
\makeatletter
\let\@oldaddcontentsline\addcontentsline
\renewcommand{\addcontentsline}[3]{}
\section*{附录}
\let\addcontentsline\@oldaddcontentsline
\makeatother

\subsection*{附件A：支撑材料说明}

\begin{table}[H]
  \centering
  \caption{支撑材料文件清单}
  \label{tab:appendix}
  \renewcommand{\arraystretch}{1.2}
  \begin{tabularx}{\textwidth}{LX}
    \toprule
    文件 & 说明 \\
    \midrule
    code/Q1/q1\_main.py & Q1主程序：Spearman相关 + K-means++聚类 + 分布拟合 \\
    code/Q1/q1\_baseline.py & Q1对比程序：Pearson相关 + 层次聚类 \\
    code/Q2/q2\_main.py & Q2主程序：Prophet分解 + 约束报童优化 \\
    code/Q2/q2\_baseline.py & Q2对比程序：ARIMA + 固定加成 \\
    code/Q3/q3\_main.py & Q3主程序：单品过滤 + 单品报童优化 \\
    code/Q3/q3\_baseline.py & Q3对比程序：固定加成策略 \\
    code/scripts/clean\_data.py & 数据清洗脚本 \\
    code/scripts/robustness\_checks.py & 稳健性检查脚本 \\
    results/QX/reports/frozen\_numbers.json & 冻结数值快照 \\
    \bottomrule
  \end{tabularx}
\end{table}

\subsection*{附件B：Q3单品级策略详表（部分）}

受篇幅限制，31个单品的完整最优策略见支撑材料中的 \texttt{q3\_m1\_sku\_policy.csv} 文件。表\ref{tab:q3sample}展示各品类代表单品的结果。

\begin{table}[H]
  \centering
  \caption{Q3部分单品最优策略}
  \label{tab:q3sample}
  \renewcommand{\arraystretch}{1.1}
  \begin{tabularx}{\textwidth}{CCCC}
    \toprule
    品类 & 入选数 & 品类最优订货量（kg） & 品类利润贡献（元） \\
    \midrule
    水生根茎类 & 3 & 17.1 & 79 \\
    花叶类 & 12 & 142.7 & 322 \\
    花菜类 & 2 & 28.2 & 131 \\
    茄类 & 3 & 39.2 & 146 \\
    辣椒类 & 6 & 132.0 & 342 \\
    食用菌 & 5 & 77.1 & 161 \\
    \bottomrule
  \end{tabularx}
\end{table}
